% % \chapter*{\Large \center Abstract}
% % Vision systems deployed in aerial, agricultural, or outdoor environments often degrade under adverse atmospheric conditions such as fog, smog, blizzards, or sandstorms. These effects induce strong distribution shifts in image statistics, reducing the reliability of vision transformers and multimodal encoders trained on clean datasets.

% % This work explores diffusion-based generative frameworks to synthetically recreate such environmental conditions, enabling vision models to adapt and generalize effectively to real-world scenes. Foundational diffusion architectures (UNET,DDPM,Conditional Diffusion) and Vision based Detection(Vit) were implemented and analyzed, followed by an evaluation of style-adaptive frameworks such as StyleShot, InstantStyle, and CSGO under extreme environmental domains.

% % Based on an existing architecure of image to image generation, an adaption was maded on InstantStyle-Plus , a diffusion framework which integrates ResNoise and IP-Adapter modules for semantic and structural consistency, and further extend it with a Contrastive Style Descriptor (CSD) Guidance Loop for perceptually aligned, distribution-shift–aware image synthesis.

% % Building upon these insights 

% % A curated dataset spanning fog, sandstorm, forest-fire, snow, smog, and underwater domains was constructed to evaluate the framework. The proposed method achieves realistic, depth-aware stylization while maintaining semantic integrity, improving robustness for agricultural, aerial, and environmental monitoring applications.
% % \chapter*{\Large \center Abstract}
% % Vision systems deployed in outdoor and aerial settings often degrade under adverse conditions such as fog, blizzard, or storm,smog, which cause significant distribution shifts in image intensity and frequency. This work explores diffusion-based generative frameworks to synthetically recreate such conditions, enabling vision models to adapt and generalize effectively to real-world environments.

% % A focused DataSet was curated having images of extreme atmospheric conditions such as SandStorm-DustStorm, ForestFire,Storm,Fog,Smog and even Deep Trences Underwater.Thee are instrumental in assessing existing ,adaptive and proposed solutions for our use case.


% % Foundational diffusion architectures, including DDPM, DDIM, and Conditional Diffusion, were studied and implemented to understand their generative behavior. Building upon these, style-adaptive frameworks such as  StyleShot, InstantStyle, and CSGO were evaluated using state of the art Vision Transformers on the extreme atmospheric images.

% % An adapted model, InstantStyle Plus, integrating ResNoise and IP-Adapter modules for semantic and structural
% % consistency, and further extend it with a Contrastive Style Descriptor (CSD) Guidance Loop for perceptually aligned, distribution-shift–aware image synthesis.


% % Finally a novel method is introduced by integrating an existing method i InstantStyle with CSD Loss, guiding the diffusion process through contrastive style supervision to achieve distribution shift–aware image synthesis. The resulting system produces realistic, style-shifted datasets that improve the robustness of vision models for agricultural, aerial, and environmental monitoring applications.

% \chapter*{\Large \centering Abstract}

% Vision systems deployed in aerial, agricultural, and outdoor environments often degrade under adverse atmospheric conditions such as fog, smog, blizzards, or sandstorms.  
% These effects cause significant distribution shifts in image statistics, reducing the reliability of vision transformers and multimodal encoders trained on clear-weather datasets.

% This work explores diffusion-based generative frameworks to synthetically recreate such environmental conditions, enabling vision models to adapt and generalize effectively to real-world scenes.  
% A focused dataset was curated containing images of extreme atmospheric domains such as fog, smog, sandstorm, forest fire, storm, and underwater to evaluate existing, adaptive, and proposed solutions under consistent conditions.

% Foundational diffusion architectures, including U–Net, DDPM, and Conditional Diffusion, were implemented to understand their generative dynamics.  
% Building on these foundations, style-adaptive frameworks such as StyleShot, CSGO, and InstantStyle were evaluated on the curated dataset using state-of-the-art vision transformers to assess their ability to reproduce physically consistent style shifts.

% An adapted model, \textbf{InstantStyle-Plus}, was developed by extending InstantStyle with ResNoise and IP-Adapter modules for semantic and structural consistency, and further enhanced through a \textbf{Contrastive Style Descriptor (CSD)} guidance loop for perceptually aligned, distribution-shift–aware image synthesis.  

% Finally, a novel diffusion system integrating \textbf{InstantStyle with CSD loss} was introduced, guiding the denoising trajectory through contrastive style supervision to achieve physically realistic, depth-aware stylization while preserving semantic fidelity.The proposed framework demonstrates superior perceptual realism and environmental consistency, producing high-quality, style-shifted datasets that enhance the robustness of vision models for agricultural, aerial, and environmental monitoring applications.

\chapter*{\Large \centering Abstract}

Vision systems deployed in aerial, agricultural, and outdoor environments often degrade under adverse atmospheric conditions such as fog, smog, blizzards, or sandstorms.  
These effects cause significant distribution shifts in image statistics, reducing the reliability of vision transformers and multimodal encoders trained on clear-weather datasets.

This work explores diffusion-based generative frameworks to synthetically recreate such environmental conditions, enabling vision models to adapt and generalize effectively to real-world scenes.  

A focused dataset was curated containing images of extreme atmospheric domains such as fog, smog, sandstorm, forest fire, storm, and underwater to evaluate existing, adaptive, and proposed solutions under consistent conditions.

Foundational diffusion architectures, including U–Net, DDPM and Conditional Diffusion were implemented to understand their generative dynamics.  
Building on these foundations, style-adaptive frameworks such as StyleShot, CSGO and InstantStyle  were evaluated on the curated dataset using state-of-the-art vision transformers  to assess their ability to reproduce physically consistent style shifts.

An existing model, \textbf{InstantStyle-Plus}, was adapted by extending InstantStyle with ResNoise  and IP-Adapter  modules for semantic and structural consistency, and further enhanced through a \textbf{Contrastive Style Descriptor (CSD)} guidance loop for perceptually aligned, distribution-shift–aware image synthesis .  

Finally, a novel diffusion system integrating \textbf{InstantStyle with CSD loss} was introduced, guiding the denoising trajectory through contrastive style supervision to achieve physically realistic, depth-aware stylization while preserving semantic fidelity.  
The proposed framework demonstrates superior perceptual realism and environmental consistency, producing high-quality, style-shifted datasets that enhance the robustness of vision models for agricultural, aerial, and environmental monitoring applications.
